\documentclass[10pt]{article}

\usepackage[utf8]{inputenc}

\usepackage[T1]{fontenc}

\usepackage{amsmath}

\usepackage{amsfonts}

\usepackage{amssymb}

\usepackage[version=4]{mhchem}

\usepackage{stmaryrd}

\usepackage{bbold}



\begin{document}

\begin{itemize}

  \item ядро П уассона для ша ра, $\sigma_{n}=$ $=n \pi^{n / 2} R^{n-1} / \Gamma(n / 2+1)$ площадь сферы $S_{n}(0, R), d S_{n}-$ элемент плоцади $S_{n}(0, R)$.

\end{itemize}



С. Пуассон [1] пришел к формуле (1) в случае $n=2$ как к интегральной форме записи суммы тригонометрич. ряда



\[

\frac{a_{0}}{2}+\sum_{k=1}^{\infty}\left(a_{k} \cos k \theta+b_{k} \sin k \theta\right) r^{k},

\]



где $a_{k}, b_{k}$ - коэффициенты Фурье функции $f(y)=f(e l \varphi)$, $(r, \theta)$ и $(1, \varphi)$ - полярные координаты соответственно точек $x=r e^{i \theta}$ и $y=e^{i \varphi}$, когда ядро Пуассона имеет вид



\[

P B_{2}(x, y)=P B_{2}\left(r e^{i \theta}, e^{i \varphi}\right)=\frac{1}{2 \pi} \frac{1-r^{2}}{1-2 r \cos (\theta-\varphi)+r^{2}}

\]



(о применениях П. и. в теории тригонометрич. рядов см. [3], а также Абеля - Пуассона метод суммирова$н и я)$.



П. и. для полупространства



\[

R_{+}^{n}=\left\{x=\left(x_{1}, \ldots, x_{n}\right) \in \mathbb{R}^{n}: x_{n}>0\right\}

\]



имеет вид



\[

u(x)=\int_{R_{0}^{n}} f(y) P R_{+}^{n}(x, y) d R_{0}^{n}(y)

\]



где



\[

R_{0}^{n}=\left\{y=\left(y_{1}, \ldots, y_{n}\right) \in \mathbb{R}^{n}: y_{n}=0\right\},

\]



$d R_{0}^{n}$ - элемент площади $R_{0}^{n}, f(y)$ - ограниченная непрерывная функция на $R_{0}^{n}$,



\[

P R_{+}^{n}(x, y)=\frac{2 \Gamma\left(\frac{n}{2}+1\right)}{n \pi^{n / 2}} \frac{x_{n}}{|x-y|^{n}}

\]



\begin{itemize}

  \item яд ро Пу ас с она для полупрос тр а н с т в а. Формулы (1) и (3) суть частные случаи формулы Грина

\end{itemize}



\[

u(x)=\int_{\Gamma} f(y) \frac{\partial G(x, y)}{\partial n_{y}} d \Gamma(y),

\]



даюгей решение задачи Дирихле для областей $D \subset \mathbb{R}^{n}$ c гладкой границей $\Gamma$ при помощи производной $d G(x, y) / d n_{y}$ функции Грина $G(x, y)$ по направлөвию внутренней нормали к $\Gamma$ в точке $y \in \Gamma$. Иногда формулу (4) также наз. П. и.



Основные свойства П. и.: 1) $u(x)$ есть гармонич. фувкция координат точки $x$; 2) П. и. дает рөшение задачи Дирихле с граничными данными $f(y)$ в классе (ограниченных) гармонич. функций, т. е. функция $u(x)$, продолженная на границу области значениями $f(y)$, непрерывна в замкнутой области. На этих свойствах основаны применения П. и. в классической математич. физике (см. [4]).



П. и., понимаемый в смысле Лебега, от суммируемой функции $f(y)$, напр. на $S_{n}(0, R)$, наз. и п т е г p aло м Пу ассона-Л е бе г а; интеграл вида



\[

u(x)=\int_{S_{n}(0, R)} \boldsymbol{P B}_{n}(x, y) d \mu(y)

\]



по произвольной конечной борелевской мере $\mu$, сосредоточенной на $S_{n}(0, R)$, наз. и в т е г p а л о м П у а сс о н а - С т и л т ь е с а. Класс $\boldsymbol{A}$ гармонич. функций $u(x)$, представимых интегралом (5), характеризуется тем, что любая функция $u(x) \in A$ есть разность двух неотрицательных гармонич. функций в $B_{n}(0, R)$. Класс функций, представимых швтегралом Пуассова - Лебега, есть правильный подкласс класса $A$, содержаций, в свою очередь, все ограниченные гармонич. функции в $B_{n}(0, R)$. Для почти всех точек $y \in S_{n}(0, R)$ по мере Лебега на $S_{n}(0, R)$ интеграл Пуассона - Стилтьеса (5) имеет угловые граничные значения, совпадающие со знаqевием производной $\mu^{\prime}(y)$ меры $\mu$ по мере Лебега. Теория интегралов Пуассона - Стилтьеса и Пуассона - Лебега строится и для случая полупространства (см. [5]).



Большую роль в теории аналитич. функций многих комплексвых переменных и в ее применениях к квантовой теории поля играют различные модификации П. и. Напр., яд ро П у а с с она для полик руга $U^{n}=\left\{z=\left(z_{1}, \ldots, z_{n}\right) \in \mathbb{C}^{n}:\left|z_{j}\right|<1, j=1, \ldots, n\right\}$ комплексвого пространства $\mathbb{C}^{n}$ получается при перемножөнии ядер (2):



\[

P U^{n}(z, \zeta)=\prod_{j=1}^{n} P B_{2}\left(z_{j}, \zeta_{j}\right)

\]



Соответствующий П. и.



\[

u(z)=\int_{T^{n}} f(\zeta) P U^{n}(z, \zeta) d T^{n}(\zeta)

\]



по остову поликруга $T^{n}=\left\{\zeta=\left(\zeta, \ldots, \zeta_{n}\right) \in \mathbb{C}^{n}:|\zeta|=\right.$,1 , $j=1, \ldots, n\}$ дает кратногармонич. функцию $u(z), z \in U^{n}$, привимающую на остове $T^{n}$ непрерывные значения $f(\zeta)$. Рассматриваются также обобщения в виде интегралов Пуассона - Лебега и Пуассона - Стилтьеса (cM. [6]).



В квантовой теории поля применяются П. и. для трубчатых областей $T C$ комплексного пространства $\mathbb{C}^{n}$ над выпуклым открытым острым конусом $C$ в пространстве $\mathbb{R}^{n}$ (с вершиной в вачале координат) вида



\[

T^{C}=\mathbb{R}^{n}+i \mathbb{C}=\left\{z=x+i y \in \mathbb{C}^{n}:\right.

\][



\textbackslash left.x=\textbackslash left(x\_\{1\}, ···, x\_\{n\}\textbackslash right) \textbackslash in \textbackslash mathbb\{R\}\^{}\{n\} ; y=\textbackslash left(y\_\{1\}, ···, y\_\{n\}\textbackslash right) \textbackslash in \textbackslash mathbb\{C\}\textbackslash right\} \textbackslash text \{. \}

]



П. и. для полуплоскости вида (3) при $n=2$ есть частный случай таких П. и. для трубчатых областей. П. и. для ограниченвых симметрич. областей пространства $\mathbb{C}^{n}$ представляется так же, как П. и. для трубчатой области в шростравстве матриц. Понимая плотность П. и. $f$ как обобщенную функцию, а сам П. и.- как свертку $f$ с ядром Пуассона, прџходят к важному понятию П. и. от обобщевных функций определенных классов (см. $[7]-[\theta])$.



Jum : [1] P o isson S. D., “J . Ecole polytechn.", 1820, t. 11, p. 295-341; 1823, t. 12, p. $404-509 ;$ [2] s c h w a r z H. A., 28; [8] В а р и Н. К., Тригонометрические ряды, М., 1961; [4] Т и х о в о в А. Н., С а м а $\mathbf{\text { в }}$ к и и А. А., Уравнения математической физики,'5 изд, М , 1977; [5] С о ло м е н ц е в Е. Д.,

итоги науки Математический анализ. Төория вероятностей. итоги науки Математический анализ. Теория вероятностећ.

Регулирование. 1962, М., 1964, с. 83-100; [6] Р у д и н У., Теория функций в поликруге, пер. с англ., М., 1974; [7] В л ад и м и р о в В. С., Обобщенные функции в математической физике. М., 1876; [8] X у а J о-к е н, Гармонический анализ функций многих комплексных переменных в классических областях, пөр. с кит., М., 1959; [9] С т е й н И., В е й с Г., Введение в гармонический анализ на евклидовых пространствах,

пер с авгл, М., 1974. ПУАССОНА МЕТОД СУММИРОВАНИЯ̆ - то же, что Абеля-Пуассона метод сумлирования.



ПУАССОНА ПРЕОБРАЗОВАНИЕ - интегральное преобразование вида



\[

f(x)=\frac{1}{\pi} \int_{-\infty}^{\infty} \frac{1}{1+(x-t)^{2}} d \alpha(t)

\]



где $\alpha(t)$ - фувкция ограниченного изменения в каждом нонечном интервале, а также преобразование



\[

f(x)=\frac{1}{\pi} \int_{-\infty}^{\infty} \frac{\varphi(t)}{1+(x-t)^{2}} d t,

\]



вытекающее из $\left({ }^{*}\right)$, если $\alpha(t)-$ абсолютно непрерывная функция. Пусть



\[

\hat{\boldsymbol{g}}(x)=-\frac{1}{\pi} \int_{0}^{\infty} u^{-2}[g(x+u)-2 g(x)+g(x-u)] d u

\]



и



\[

\begin{gathered}

T_{t} g(x)=\sum_{k=0}^{\infty}(-1)^{k} \frac{t^{2 k}}{(2 k) !} g^{(2 k)}(x)+ \\

+\sum_{k=0}^{\infty}(-1)^{k} \frac{t^{3 k+1}}{(2 k+1) !} \hat{g}^{(2 k)}(x) .

\end{gathered}

\]





\end{document}